% Standalone paragraph. Paste into your own thesis, or \input it.
% Not wired into any chapter's numbering -- drop it wherever the "which features are actually
% reproducible" argument is needed. Companion to whole_tree_gap.tex: that file argues the field
% measures the wrong place in the tree, this one argues it does not agree on how to measure what
% it does look at.

% HEADLINE: one of three sibling subsections for the State of the Art chapter, sitting under
% \section{The gap} (state_of_the_art/08_the_gap.tex): where the tree is measured
% (whole_tree_gap), how it is measured (reproducible_features), in whom it is measured
% (healthy_cohort_2). Promote to \section if they are to stand as sections of their own.
\subsection{How the tree is measured}
\label{sec:gap-how}

No single cohort is large enough to settle whether a geometric feature predicts disease; that
power comes only from comparing findings across studies, and comparison is only valid once every
study measured the same quantity.
Two coronary features are standardized, and for the same reason: each has an immediate procedural
use that forces one number.
Segmentation work now treats the centerline, not the lumen mask, as the shared unit of comparison
\cite{shit2021cldice, zhang2024adehtl, nannini2024characterization}, and bifurcation angle has a
fixed definition because a stent has to be placed the same way every time, codified across eighteen
consensus documents from the European Bifurcation Club \cite{burzotta2024ebc}.

Curvature and tortuosity have no such settlement, and not for the same reason.

Curvature is mathematically well-posed -- the field's own recent review gives one formula,
$\kappa = |\mathrm{d}T/\mathrm{d}s|$, equal to $1/R$ for the osculating circle -- but is not applied
consistently: computational studies routinely collapse a curved segment to a single bend read off
that formula, ``neglecting the spatiality and continuity of this characteristic''
\cite{shen2026geometry}.
The problem is discretization, not definition.

Tortuosity's problem is more fundamental than that: its dominant metric, the tortuosity index, compares path
length to straight-line endpoint distance and, in the review's own words, ``cannot capture the
spatial information'' in the first place \cite{shen2026geometry}.
Clinical practice reduces it further still, to a bend count or a 2D C-/S-shape projection.
The formula itself, not just its use, is contested.

Pointed at the same vessels, the metrics disagree outright.
Kashyap et al.\ tested every tortuosity metric then in use against low wall shear stress in a single
127-patient cohort: average absolute curvature was the strongest correlate ($p<0.001$), while the
tortuosity index --- the most-used measure in the literature --- showed none at all ($p=0.86$)
\cite{kashyap2022tortuosity}.
Zebi\'c Mihi\'c et al.\ report the same split from the clinic: a continuous tortuosity index caught
ischaemia in 160 patients that the standard bend-angle count missed entirely
\cite{zebicmihic2024index}.
Even Nannini et al.'s automated pipeline declines to choose, reporting three incompatible tortuosity
metrics side by side rather than committing to one \cite{nannini2024characterization,
zebicmihic2023tortuosity}.

The disagreement reaches outcomes, too.
Two angiographic cohorts both found tortuosity protective against significant disease, but did not
measure the same thing to get there: one counted two consecutive $180^\circ$ turns and called
12.45\% of 1221 patients tortuous \cite{groves2009tortuosity}, the other counted three or more bends
past $45^\circ$ and called 39.1\% of 1010 \cite{li2011tortuosity}.
Meanwhile the same review's evidence table reports continuous curvature measures pointing the other
way at the segment level \cite{shen2026geometry}.
Sex confounds even that split --- tortuosity is markedly more common in women
\cite{li2011tortuosity} --- and neither cohort's headline result was stratified by it.
Whether tortuosity helps or harms a patient looks less like a biological fact than an artifact of
which formula, scale, and sex mix a given study happened to use.

One of those results was obtained in the population this thesis measures, and that decides it.
Kashyap et al.\ selected average absolute curvature in 127 patients \emph{without} coronary artery
disease, against a criterion --- the vessel area below 0.4~Pa of time-averaged shear --- that exists
in an artery with no plaque in it \cite{kashyap2022tortuosity}.
Every rival criterion needs disease to already be present to rank one metric above another.
This thesis therefore computes average absolute curvature along the centerline, point by point
rather than collapsed to a single bend, and reports a continuous index rather than a bend count
wherever a tortuosity-style quantity is reported at all \cite{zebicmihic2024index}.
Two limits come with that choice.
Kashyap measured at the left main and its branches, so applying the metric across the whole tree
extrapolates beyond where it was validated; and Zhang et al.\ find the same quantity separating
stenosed from non-stenosed trees, but cross-sectionally and in symptomatic patients, which is why
they call for longitudinal studies starting from the healthy state rather than resting on the
comparison \cite{zhang2024curvature}.

Measuring an unselected population changes what validation can even mean.
There is no stenosis to discriminate against, so the metric cannot be judged by how well it
separates diseased trees from healthy ones; it is judged instead by whether the same person scanned
twice yields the same number, and only then against what happens to that person --- the measurement
floor before the association, since two accepted techniques for measuring bifurcation angle have
been shown to disagree by almost the entire width of the clinical gradient they were meant to
resolve \cite{givehchi2018phantom}.
Sex stratification stops being optional for the same reason.
Tortuosity is markedly more common in women \cite{li2011tortuosity}, and the one published normative
distribution over disease-free vessels found men and women differing even after adjustment for lumen
diameter \cite{schultz2023ess}.
Pooled into a single reference distribution, a healthy woman's tree would be measured against a
population it does not belong to.

A feature is reproducible only once its definition is fixed before the analysis starts, not
recovered afterwards to explain why one paper's finding is the sign flip of another's.
